\subsection{DCGAN}

Generative Adversarial Networks (GANs) are a class of machine learning frameworks introduced by Goodfellow et al. in 2014 \cite{goodfellow2014}. GANs have revolutionized the field of generative modeling, allowing for the creation of highly realistic synthetic data. This section will explore the fundamental concepts, structure, and applications of Deep Convolutional Generative Adversarial Networks (DCGANs).

\subsubsection{Introduction to DCGANs}

DCGANs consist of two neural networks: the generator and the discriminator, which compete to minimize their respective loss functions. The generator aims to create data that is indistinguishable from real data, while the discriminator attempts to differentiate between real and generated data. This adversarial process continues until the generator produces data that the discriminator cannot reliably distinguish from real data. DCGANs leverage convolutional neural networks (CNNs) to improve the quality of generated images.

\begin{figure}[H]
    \centering
    \includegraphics[width=0.6\textwidth]{assets/gan_structure.png}
    \caption{Structure of a Generative Adversarial Network, picture from \cite{kdnuggets2017}}
    \label{fig:gan_structure}
\end{figure}

\subsubsection{Components of DCGANs}

The two primary components of DCGANs are:

\paragraph{Generator}
\hfill \break
The generator, \( G \), is a convolutional neural network that takes random noise, \( z \), from a latent space and transforms it into a synthetic data sample, \( G(z) \). The objective of the generator is to produce data that is as close as possible to the real data distribution. The generator's loss function is designed to minimize the likelihood of the discriminator correctly identifying the generated data as fake.
\\
The actual structure of the generator, inspired by the standard DCGAN presented in \cite{radford2016}, is reported in table \ref{tab:dcgan generator}.

\begin{table}[H]
    \centering
    \begin{tabular}{l c r}
    \textbf{Layer} & \textbf{Output shape} & \textbf{Parameters} \\
    \hline
    Dense & \( 1152 \) & 24{,}192 \\
    Reshape & \( 3 \times 3 \times 128 \) & -- \\
    Transposed convolution & \( 6 \times 6 \times 128 \) & 262{,}272 \\
    LeakyReLU & \( 6 \times 6 \times 128 \) & -- \\
    Transposed convolution & \( 14 \times 14 \times 256 \) & 524{,}544 \\
    LeakyReLU & \( 14 \times 14 \times 256 \) & -- \\
    Transposed convolution & \( 28 \times 28 \times 512 \) & 2{,}097{,}664 \\
    LeakyReLU & \( 28 \times 28 \times 512 \) & -- \\
    Convolution & \( 28 \times 28 \times 1 \) & 12{,}801 \\
    \hline
    \textbf{Total} & & \textbf{2{,}921{,}473} \\
    \end{tabular}
    \caption{Structure of the DCGAN generator, shown for a latent dimension of 20. With a latent dimension of 30 only the first dense layer changes, and the total becomes 2{,}932{,}993 parameters.}
    \label{tab:dcgan generator}
\end{table}

\paragraph{Discriminator}
\hfill \break
The discriminator, \( D \), is another convolutional neural network that evaluates data and assigns a probability that a given sample comes from the real data distribution rather than the generator. The discriminator's loss function aims to maximize the likelihood of correctly classifying real and fake data.
\\
The actual structure of the discriminator, inspired by the standard DCGAN presented in \cite{radford2016}, is reported in table \ref{tab:dcgan discriminator}.

\begin{table}[H]
    \centering
    \begin{tabular}{l c r}
    \textbf{Layer} & \textbf{Output shape} & \textbf{Parameters} \\
    \hline
    Convolution & \( 14 \times 14 \times 64 \) & 1{,}088 \\
    LeakyReLU & \( 14 \times 14 \times 64 \) & -- \\
    Convolution & \( 7 \times 7 \times 128 \) & 131{,}200 \\
    LeakyReLU & \( 7 \times 7 \times 128 \) & -- \\
    Convolution & \( 4 \times 4 \times 128 \) & 262{,}272 \\
    LeakyReLU & \( 4 \times 4 \times 128 \) & -- \\
    Flatten & \( 2048 \) & -- \\
    Dropout & \( 2048 \) & -- \\
    Dense & \( 1 \) & 2{,}049 \\
    \hline
    \textbf{Total} & & \textbf{396{,}609} \\
    \end{tabular}
    \caption{Structure of the DCGAN discriminator}
    \label{tab:dcgan discriminator}
\end{table}

\subsubsection{Training Process}

The training process for DCGANs involves optimizing the following value function, \( V(G, D) \), which is a minimax game between \( G \) and \( D \):
\begin{equation}
\min_G \max_D V(G, D) = \mathbb{E}_{\mathbf{x} \sim p_\text{data}(\mathbf{x})} [\log D(\mathbf{x})] + \mathbb{E}_{\mathbf{z} \sim p_\mathbf{z}(\mathbf{z})} [\log (1 - D(G(\mathbf{z})))]
\end{equation}
The training of DCGANs involves alternating between updating the discriminator and the generator. The discriminator is trained to maximize its classification accuracy, while the generator is trained to deceive the discriminator. 

One detail of our implementation is worth mentioning: the labels given to the discriminator are perturbed by a small amount of uniform noise, drawn in \( [0, 0.05] \) and applied so that each class moves towards the other. This keeps the discriminator from saturating early in training, which would leave the generator without a useful gradient. The direction matters: adding the same positive noise to both classes, as the widely copied reference implementation does, pushes the positive targets past 1 and rewards overshooting instead.
\\
The procedure can be summarized in the following steps:

\begin{enumerate}
    \item Sample a batch of real data and a batch of random noise.
    \item Update the discriminator by minimizing its classification loss on both real and generated data, which corresponds to maximizing \( V(G,D) \).
    \item Sample a new batch of random noise.
    \item Update the generator by minimizing its loss function, which depends on the discriminator's performance.
\end{enumerate}

\subsubsection{Training Configuration}

Both networks are trained on the 54,000 images the VAE trains on, which is the training split with the same tenth held out, so that the two families see the same data and the test split takes no part in either. Training runs for 50 epochs with mini-batches of 64, one Adam optimizer for each network at a learning rate of 0.0002 with \( \beta_1 = 0.5 \), a binary cross-entropy loss, and two generator updates per discriminator update. These are the settings \cite{bora2017} give for their DCGAN in section 5.2. The learning rate and the momentum term are those of \cite{radford2016}, whose architecture that paper adopts and which singles out a high \( \beta_1 \) as a cause of unstable training; the mini-batch size of 64 is the reference paper's own choice. Two versions were trained, one with a latent vector of dimension 20 and one of dimension 30.
\\
A generator is saved every five epochs from the twentieth, giving seven candidates per run, and the one used in section 6 is chosen among them by a rule fixed before training began: the lowest representation error over 32 images held out of training. That rule is not a formality, because the error does not fall monotonically with the epoch: it spans a factor of 1.28 at \( k = 20 \) and 1.42 at \( k = 30 \) across the saved candidates, and in both runs the final epoch was about 11\% worse than the one the rule selected. Keeping whatever the last epoch produced, which is the usual practice when a model has no validation loss, would have given a visibly weaker prior.


